% Auxiliary Hodge input for the higher-range local restriction construction.
% This fragment proves the Hodge assertion only, not the local geometry.
\paragraph{Auxiliary product Hodge calculation for quartics containing lines.}
Let $s\geq1$. We choose smooth complex quartics $W_1,\ldots,W_s$, each
containing a specified line $D_i$, at a jointly Hodge-generic point of
the product of the corresponding quartic families. Write
$H_i=c_1(\mathcal O_{W_i}(1))$ and
$T_i=(\mathbb QH_i+\mathbb QD_i)^\perp\subset H^2(W_i,\mathbb Q)$.
We can and do require
\begin{equation}\label{eq:product-K3-generic-data}
 \begin{gathered}
 \operatorname{NS}(W_i)=\mathbb ZH_i\oplus\mathbb ZD_i,
 \qquad
 \begin{pmatrix}H_i^2&H_iD_i\\H_iD_i&D_i^2\end{pmatrix}
       =\begin{pmatrix}4&1\\1&-2\end{pmatrix},\\
 \operatorname{End}_{\mathrm{Hdg}}(T_i)=\mathbb Q,
 \qquad
 \operatorname{Hom}_{\mathrm{Hdg}}(T_i,T_j)=0\quad(i\ne j).
 \end{gathered}
\end{equation}
Let $C$ be the Fermat plane quartic, set
$S=(\prod_{i=1}^sW_i)\times C$, and consider the indicated Kunneth
summand
\[
 K=\left(\bigotimes_{i=1}^sH^2(W_i,\mathbb Q)\right)
                       \otimes H^1(C,\mathbb Q)
       \subset H^{2s+1}(S,\mathbb Q).
\]
Then its largest rational sub-Hodge structure of level at most one is
\begin{equation}\label{eq:product-K3-level-one}
 K_{\leq1}=
 \left(\bigotimes_{i=1}^s\operatorname{NS}(W_i)_{\mathbb Q}\right)
                       \otimes H^1(C,\mathbb Q).
\end{equation}
In particular, for every smooth projective complex curve $P$ and every
$z\in\mathrm{CH}^{s+1}(S\times P)$, the component of its rational
cycle class in $K\otimes H^1(P,\mathbb Q)$ belongs to
$K_{\leq1}\otimes H^1(P,\mathbb Q)$. The parameter curve is arbitrary;
it is not required to be defined over the field of definition of the
quartics. We assert \eqref{eq:product-K3-level-one} for the specified
Kunneth summand $K$, not for all of $H^{2s+1}(S,\mathbb Q)$.

\emph{Existence of the chosen tuple.}\par
For a fixed line $D_0\subset\mathbb P^3$, let $B$ be the smooth open
subset of $\mathbb P H^0(\mathcal I_{D_0}(4))$. Theorem~1.1 and
Corollary~1.3(c) of \citet[p.~2 of the opened author version]
{frBrevikNollet2011} give $\operatorname{Pic}(W)=\mathbb ZH\oplus
\mathbb ZD_0$ for a very general point of $B$. Their hypotheses hold
because $\mathcal I_{D_0}(3)$ is globally generated and
$H^0(\mathcal I_{D_0}(2))\ne0$. Adjunction gives the intersection matrix
in \eqref{eq:product-K3-generic-data}.

We impose the additional Hodge conditions explicitly. By Local Torelli
and the line-bundle period description in
\citet[Proposition~6.2.8 and Section~6.2.4, pp.~107--108 of the
opened author version]{frHuybrechtsK3Book}, locally marked deformations
preserving $H,D_0$ fill an open subset of the $18$-dimensional period
quadric in $(\mathbb QH+\mathbb QD_0)^\perp_{\mathbb C}$.
They are locally realized in $B$: very ampleness of $H$ is open;
Riemann--Roch and $H\cdot D_0=1$ make the effective class $D_0$ a line;
and a projective coordinate change returns that line to the fixed line.
Thus the product family has independently varying local periods.

For a nonscalar rational endomorphism $a$ of a marked transcendental
space, the Hodge-endomorphism condition forces the period line to be
an eigenline of $a$. The finite union of its eigenspaces cannot contain
the irreducible nondegenerate period quadric, which spans the ambient
projective space. This condition therefore defines a proper locus.
After imposing $\operatorname{End}_{\mathrm{Hdg}}(T_i)=\mathbb Q$,
each $T_i$ is simple: a nontrivial direct summand would give a
nontrivial rational Hodge projector. A nonzero Hodge map between two
$T_i$ would consequently be an isomorphism. For each rational linear
isomorphism $a:T_i\to T_j$, its Hodge locus forces
$a\omega_i\in\mathbb C\omega_j$. The independent variation of the
two period lines makes this a proper locus in the product family.
There are only countably many such rational maps. Applied to the
weight-zero variations $\operatorname{End}(T_i)$ and
$\operatorname{Hom}(T_i,T_j)$, the Hodge-locus theorem
\citep[Theorem~1.1 and Corollary~1.3, p.~2]{frCDK1994}
makes the loci, with their monodromy translates, countable unions of
closed algebraic subsets. Avoiding these loci and the extra Picard
loci gives \eqref{eq:product-K3-generic-data}. This specifies a jointly
Hodge-generic complex point; it does not assert that every embedding
of a coefficientwise algebraically independent parameter field is
Hodge generic.

\emph{Independence of the orthogonal factors.}\par
By \citet[Theorem~2.2.1 and Example~2.2.3, p.~211]{frZarhin1983},
each $T_i$ has rank $20$ and $\operatorname{Hg}(T_i)=\operatorname{SO}(T_i)$.
Let $G=\operatorname{Hg}(\bigoplus_iT_i)$ and
$\mathfrak g=\operatorname{Lie}(G)$. Its projections onto the
absolutely simple Lie algebras
$\mathfrak g_i=\mathfrak{so}(T_i)$ are surjective.
The radical of $\mathfrak g$ maps to a solvable ideal in every
$\mathfrak g_i$, hence is zero. Decompose this semisimple rational
Lie algebra into its simple ideals. Each projection to
$\mathfrak g_i$ is an isomorphism on exactly one of these ideals and
zero on the others. Consequently, if
$\mathfrak g\ne\bigoplus_i\mathfrak g_i$, one simple ideal maps
isomorphically onto two different factors. The projection to that
pair is the graph of a rational Lie-algebra isomorphism
$\varphi:\mathfrak g_i\simeq\mathfrak g_j$.

We record why this would identify the two standard representations;
this is the point at which the rank matters. Over $\overline{\mathbb Q}$,
an automorphism of $\mathfrak{so}_{20}$ integrates to the simply
connected group, and, modulo an inner automorphism, acts by a
symmetry of its $D_{10}$ diagram
\citep[20.33, p.~108, and 21.4, p.~112]{frMilneReductive2018}.
The simple roots are
$\varepsilon_1-\varepsilon_2,\ldots,
\varepsilon_9-\varepsilon_{10},\varepsilon_9+\varepsilon_{10}$
\citep[18.26--18.27, pp.~91--92]{frMilneReductive2018}.
The only nontrivial diagram symmetry swaps its two short terminal
branches and fixes the first node. It therefore fixes the highest
weight $\varepsilon_1$ of the standard representation, whose weights
are $\{\pm\varepsilon_j\}_{j=1}^{10}$. Highest-weight uniqueness
\citep[20.3--20.4, pp.~102--103]{frMilneReductive2018}
shows that twisting the standard representation by any such
automorphism preserves its isomorphism class.

It follows that the intertwiner space
\[
 \{u\in\operatorname{Hom}_{\mathbb Q}(T_i,T_j):
          uX=\varphi(X)u\text{ for every }X\in\mathfrak g_i\}
\]
is nonzero after scalar extension. These are rational linear equations,
so the rational space itself is nonzero. A nonzero intertwiner is an
isomorphism, since both standard representations are absolutely
irreducible. It commutes with $\mathfrak g$ on the pair, hence with
the connected group $G$. Because $T_i,T_j$ have the same weight, it
is a rational Hodge isomorphism, contrary to
\eqref{eq:product-K3-generic-data}. We conclude that
\[
       \operatorname{Hg}\left(\bigoplus_{i=1}^sT_i\right)
                          =\prod_{i=1}^s\operatorname{SO}(T_i).
\]

\emph{The CM factor and the levels of the tensor constituents.}\par
By \citet[Lemma~2, p.~106, and Example~3, pp.~106--107]
{frBeauville2014}, $\operatorname{Jac}(C)$ is isogenous to $E^3$,
where $E=\mathbb C/\mathbb Z[i]$. Put $V=H^1(E,\mathbb Q)$ and
$U=\operatorname{Res}^1_{\mathbb Q(i)/\mathbb Q}\mathbb G_m$.
Then $H^1(C,\mathbb Q)\simeq V^{\oplus3}$ and
$\operatorname{Hg}(V)=U$. The joint Hodge group is a connected
subdirect product of $\prod_i\operatorname{SO}(T_i)$ and $U$.
Its derived group projects onto the full semisimple product and
trivially onto $U$, so it contains the semisimple product. Surjectivity
onto $U$ now gives
\[
 \operatorname{Hg}\left(\bigoplus_iT_i\oplus V\right)
             =\left(\prod_i\operatorname{SO}(T_i)\right)\times U.
\]

For a nonempty subset $I\subset\{1,\ldots,s\}$, put
$T_I=\bigotimes_{i\in I}T_i$ and $t=|I|$.
The representation $T_I$ of the product of the corresponding
orthogonal groups is absolutely irreducible. Over $\mathbb Q(i)$,
$T_I\otimes V$ is the sum of its twists by the two opposite
nontrivial characters of $U$. These are nonisomorphic absolutely
irreducible representations, exchanged by complex conjugation.
Therefore $T_I\otimes V$ is a simple rational Hodge structure.
After the Tate twist supplied by the remaining $s-t$ Neron--Severi
factors, it has a nonzero component of type
\[
                       (s+t+1,s-t),
\]
and hence has level $2t+1\geq3$. Each remaining Neron--Severi
factor has rank two, and $H^1(C)=V^{\oplus3}$; these give
multiplicities, not further simple constituents of smaller level.

Expand $K$ using $H^2(W_i)=\operatorname{NS}(W_i)_{\mathbb Q}\oplus T_i$.
The term with no transcendental factor is precisely the right-hand
side of \eqref{eq:product-K3-level-one}, and has types
$(s+1,s),(s,s+1)$. Every other term is a direct sum of the simple
constituents just described. Semisimplicity of polarizable rational
Hodge structures shows that none of these other terms has a nonzero
sub-Hodge structure of level at most one. This proves
\eqref{eq:product-K3-level-one}.

Finally, the selected Kunneth component of an algebraic cycle class
$\operatorname{cl}(z)$ is a rational Hodge class of type $(s+1,s+1)$.
Poincare duality on $P$ identifies it with a Hodge morphism
\[
                 H^1(P,\mathbb Q)(-s)\longrightarrow K.
\]
Its source has level at most one. Its image therefore lies in
$K_{\leq1}$, proving the stated consequence for every parameter $P$.
